\documentclass[11pt,a4paper,roman]{moderncv}
\moderncvstyle{banking}
\moderncvcolor{purple}
\nopagenumbers{}
\usepackage[utf8]{inputenc}
\usepackage[T1]{fontenc}
\usepackage{lmodern}
\usepackage[scale=0.96,top=0.48cm,bottom=0.48cm,left=0.72cm,right=0.72cm]{geometry}
\renewcommand{\baselinestretch}{0.92}
\setlength{\hintscolumnwidth}{2.20cm}
\setlength{\footskip}{18pt}

\name{Wahid}{SAMER}
\address{Cergy (95), France}{}{}
\mobile{07 53 10 95 74}
\email{samer.ahmed.wahid@gmail.com}
\social[linkedin]{Wahid SAMER}
\social[github]{wahidrt}
\title{Alternance M2 -- Études actuarielles / Modélisation financière}

\begin{document}
\makecvtitle
\vspace{-3.4em}
\begin{center}
\parbox{0.97\textwidth}{\centering\small\itshape
Étudiant en \textbf{Master Mathématiques Appliquées à l'Ingénierie Financière}, avec une formation en \textbf{modélisation stochastique, mathématiques financières, simulation et méthodes numériques}. Expérience académique en \textbf{Python, Excel/VBA et C/C++}. Souhaite appliquer ces compétences aux problématiques de \textbf{rentabilité, épargne/retraite et risque}. Recherche d'une \textbf{alternance M2 de 12 mois à partir du 21 septembre 2026}.}
\end{center}
\vspace{0.12em}

\section{Formation}
\cventry{2025 -- 2027}{Master Mathématiques Appliquées à l'Ingénierie Financière}{CY Cergy Paris Université / CY Tech}{Cergy}{}{\textbf{M1 :} probabilités, processus stochastiques, simulation stochastique, EDP, méthodes numériques avancées, optimisation, évaluation d'actifs contingents, gestion de portefeuille, Python/VBA.\\
\textbf{M2 2026--2027 :} stochastic calculus, model calibration and simulation, séries temporelles, machine learning avec Python, portfolio management, fixed income.}
\cventry{2021 -- 2025}{Licence Mathématiques et Applications}{Université d'Évry Paris-Saclay}{Évry}{}{Probabilités, mesure et intégration, analyse, algèbre, topologie, analyse numérique et programmation.}

\section{Projets ciblés Modélisation, Risque \& Données}
\cventry{2025}{Excel / VBA}{Pricer d'options vanilles et analyse de sensibilités}{Projet M1}{}{Développement d'un outil de valorisation Call/Put sous Black--Scholes ; calcul des Greeks et scénarios de stress sur la volatilité pour mesurer l'impact sur les prix.}
\cventry{2026}{Python / \LaTeX}{Modélisation stochastique -- Black--Scholes \& Heston}{Mémoire M1}{}{Étude d'Itô et Feynman--Kac, simulation et méthodes numériques ; analyse des hypothèses de modèle et comparaison avec une volatilité stochastique.}
\cventry{2026}{Python}{Modélisation statistique et risque d'un portefeuille multi-actifs}{Projet M1 -- équipe de 4}{}{Traitement de données financières, cointégration, VECM, filtre de Kalman et Ornstein--Uhlenbeck ; suivi du P\&L, drawdown et métriques de performance/risque.}

\section{Compétences}
\cvitem{Modélisation}{Probabilités, processus stochastiques, simulation, méthodes numériques, optimisation, analyse de sensibilités.}
\cvitem{Finance / Risque}{Mathématiques financières, produits dérivés, pricing, Greeks, gestion de portefeuille, stress tests.}
\cvitem{Programmation}{Excel/VBA ; Python (NumPy, Pandas, SciPy, Matplotlib) ; C/C++, SQL ; Git/GitHub, \LaTeX.}
\cvitem{Langues}{Français : courant ; Arabe : langue maternelle ; Anglais : intermédiaire.}

\section{Expérience \& Engagement}
\cventry{Mai -- juin 2025}{Professeur de mathématiques -- stagiaire}{Lycée Pierre Mendès France}{Vitrolles}{}{Préparation et transmission de notions de Terminale ; rigueur, pédagogie, capacité d'abstraction et esprit de synthèse.}
\cventry{2019}{Technicien aéronautique -- stagiaire}{G1 Aviation}{Gap}{}{Assemblage/maintenance dans un environnement exigeant en précision, procédures et sécurité.}
\cvitem{Associatif}{Bénévole aux \textbf{Restaurants du Cœur} pendant près de 3 ans : solidarité, accueil et travail d'équipe.}

\end{document}
